% Compile with pdfLaTeX.

\documentclass[14pt,a4paper,oneside]{extreport}

% Fonts and language
\usepackage[T1]{fontenc}
\usepackage[utf8]{inputenc}
\usepackage{mathptmx}
\usepackage[english]{babel}

% Page layout
\usepackage[
  left=25mm, right=22mm,
  top=40mm, bottom=25mm,
  headheight=19pt, headsep=24pt
]{geometry}
\usepackage{setspace}
\setstretch{1.5}
\usepackage{microtype}
\usepackage{indentfirst}

% Tables and graphics
\usepackage{booktabs,longtable,array,calc}
\usepackage{graphicx}
\usepackage{xcolor}
\usepackage{etoolbox}
\usepackage{needspace}
\usepackage{fancyhdr}

% Citations: separate brackets, without range compression.
\usepackage[numbers,sort]{natbib}
\setcitestyle{citesep={], [}}
\usepackage{xurl}
\usepackage[
  unicode,
  hidelinks,
  bookmarksnumbered=true
]{hyperref}
\setlength{\emergencystretch}{3em}
\setlength{\parindent}{10mm}
\setlength{\parskip}{0pt}
\widowpenalty=10000
\clubpenalty=10000
\brokenpenalty=10000
\raggedbottom

% Chapter heading spacing
\makeatletter
\patchcmd{\@makechapterhead}
  {\vspace*{50\p@}}
  {\vspace*{110\p@}}
  {}{}
\patchcmd{\@makeschapterhead}
  {\vspace*{50\p@}}
  {\vspace*{110\p@}}
  {}{}
\makeatother
\providecommand{\tightlist}{%
  \setlength{\itemsep}{0pt}%
  \setlength{\parskip}{0pt}%
}
\providecommand{\passthrough}[1]{#1}
\BeforeBeginEnvironment{longtable}{\Needspace{100mm}}
\AtBeginEnvironment{longtable}{%
  \small%
  \setstretch{1.0}%
  \setlength{\tabcolsep}{5pt}%
}
\setlength{\LTcapwidth}{\textwidth}
\setcounter{tocdepth}{2}

% Running headers
\pagestyle{fancy}
\fancyhf{}
\fancyhead[L]{%
  \small\bfseries%
  \parbox[b]{0.89\textwidth}{\nouppercase{\leftmark}}%
}
\fancyhead[R]{\small\bfseries\thepage}
\renewcommand{\headrulewidth}{0.4pt}
\renewcommand{\sectionmark}[1]{\markboth{\thesection\quad #1}{}}
\renewcommand{\chaptermark}[1]{\markboth{#1}{}}
\fancypagestyle{plain}{%
  \fancyhf{}%
  \renewcommand{\headrulewidth}{0pt}%
}

\AtBeginDocument{%
  \selectlanguage{english}%
  \renewcommand{\bibname}{References}%
}
\renewcommand{\bibname}{References}
\hypersetup{
  pdftitle={Development and Experimental Evaluation of a Self-Adaptive Computational Task Scheduling
    System for Cross-Platform Flutter Applications},
  pdfauthor={Amir Aminov}
}
\begin{document}

\pagenumbering{arabic}
\setcounter{page}{2}
\tableofcontents
\clearpage
\chapter{Introduction}
\label{introduction}

\section{Research context and motivation}
\label{research-context-and-motivation}

Native Flutter applications execute Dart logic in isolates: concurrent execution units with
their own event loops and isolated mutable application state, communicating through messages.
The main isolate processes Dart-side input callbacks and prepares the user interface for
rendering \citep{flutter-isolates}. In the standard Android and Windows embedders, its
Dart code runs on the platform main thread by default. UI and platform threads were merged
by default on Android and iOS in Flutter 3.29, and on Windows and macOS in Flutter 3.35
\citep{flutter-thread-merge}. Consequently, a long synchronous computation can delay Dart
event handling, frame processing and platform work scheduled on the same thread. Flutter
recommends background isolates for suitable CPU-intensive operations and distinguishes
one-off computations from repeated work in long-lived workers. Parsing structured data, compression,
hashing and local indexing are examples for which execution placement can matter.

The practical objective is to complete useful computations while keeping the interface
responsive. Frame smoothness and response to input are related outcomes with different
measurements. Flutter reports frame build and raster durations and the wait between receiving
vsync and starting frame processing \citep{flutter-frame}. A task can delay that start even
when subsequent stages remain short. An empirical Android study by Lin et al.~identifies
several causes of poor responsiveness: rendering work, I/O and scheduling delays
\citep{lin2024}. Execution placement must therefore be evaluated against both stage duration
and delays surrounding frame processing.

Background execution also has computational and memory costs. A new isolate must be started,
inputs made available and results returned. A persistent worker can retain initialized data
and reuse it, but this state occupies memory between tasks, and submitted work can wait in its
queue. Ordinary messages may copy mutable data, whereas \texttt{Isolate.run} returns its
result through a terminating isolate without copying it \citep{dart-run, dart-send}. The
useful execution context can consequently depend on computation, representation, repetition,
device and current application conditions.

\section{Problem statement and research gap}
\label{problem-statement-and-research-gap}

This thesis considers three alternatives for an explicitly submitted CPU-bound Dart task:
execution on the main isolate, execution in a short-lived isolate and execution in a
persistent worker. Direct execution avoids isolate dispatch and communication, but occupies
the UI event loop. Short-lived execution creates a separate isolate for the computation.
Persistent execution reuses an existing worker. Selecting among these alternatives requires
considering both task performance and the UI workload with which the task competes.

Runtime execution choice is established. Wagle et al.~(2019) select inline or asynchronous
execution using observed task costs and adapt the threshold in an HPX-based runtime
\citep{wagle2019}. Furmento et al.~(2025) adapt task granularity in StarPU using runtime
performance models and available parallelism \citep{furmento2025}. iAware changes the timing
of service work according to interaction and system state \citep{iaware2023}. These approaches
provide precedents for adaptation with different execution mechanisms and performance
objectives.

The proposed investigation concerns Dart isolate lifecycle, input and result transfer, worker
availability and Flutter frame telemetry. Existing libraries provide substantial execution
infrastructure. The unresolved question is when selecting an execution context helps and
whether recent frame information improves a decision already informed by task characteristics,
execution history and queue state. The study will test this relationship under native Flutter
UI workloads, including conditions in which a fixed strategy is sufficient or adaptation
loses.

\section{Research object, aim and objectives}
\label{research-object-aim-and-objectives}

The object of research is the execution of explicitly registered CPU-bound Dart tasks in
native Flutter applications under UI workloads. The investigation focuses on the relationship
between execution context, task and data characteristics, runtime conditions, and the
resulting trade-off between frame performance and task cost.

The aim is to develop and experimentally evaluate an application-level execution selector and
to determine the conditions under which its decisions are useful relative to fixed execution
strategies and hybrid selection policies. Usefulness concerns a trade-off across frame
performance, task completion and
execution cost: a better frame statistic alone will not establish an overall improvement when
required work is omitted or task latency becomes unacceptable. The methodology will define the
operational comparison criteria. The objectives are to:

\begin{enumerate}
  \def\labelenumi{\arabic{enumi}.}
  \tightlist
  \item
    Characterize the three execution modes across computation, data-transfer and submission
  patterns.
  \item
    Implement an open-source Flutter package with an integrable task interface and an
  explainable runtime selection policy.
  \item
    Compare the policy with fixed execution strategies, a calibrated hybrid policy with fixed
  thresholds, and the same policy with UI feedback disabled under equivalent workloads.
  \item
    Assess frame-processing delay and stage duration, task latency, throughput, memory use, and
  monitoring and scheduling overhead.
  \item
    Quantify calibration cost and evaluate policy transfer across device--OS configurations
  using reproducible workloads and separate calibration and evaluation runs.
\end{enumerate}

\section{Research questions}
\label{research-questions}

The investigation addresses three questions:

\textbf{RQ1.} How do the costs of main-isolate execution, short-lived isolates, and persistent
workers change with computation size, transferred data, task frequency, and device
performance?

\textbf{RQ2.} How does a runtime execution-selection policy affect frame-processing delay,
frame-budget exceedances, task latency and scheduling overhead relative to fixed execution
strategies and hybrid selection baselines,
and under which workload and runtime conditions are the trade-offs favorable?

\textbf{RQ3.} How do calibration budget and transfer to previously unseen device--OS
configurations affect frame performance, task latency and scheduling overhead?

RQ1 establishes the cost landscape. RQ2 tests benefits and regressions against both fixed
execution modes and a calibrated hybrid policy whose decision thresholds remain fixed during
evaluation. The latter tests whether runtime adaptation improves on a rule fitted in advance.
The contribution of UI feedback is evaluated separately by disabling only the proposed
policy's UI inputs. Task-cost estimation, execution history and queue-state inputs remain
enabled. Both variants start from matched estimates and learning budgets. Their later
histories can differ because their decisions differ.

Frame-stage budget exceedances will be treated as indicators of expensive frame processing,
with the display refresh interval as reference. Waiting before frame processing will be
examined separately using \texttt{vsyncOverhead} and frame-start intervals during continuous
animation. These observations do not directly measure physical input-to-display latency
\citep{flutter-frame, flutter-vsync, flutter-build}. The methodology and pilot will establish
operational metrics and confirm that they detect deliberately introduced CPU blocking before
evaluating the selector. RQ3 examines calibration effort and transfer between deployment
configurations, each defined by a device and its operating system. Android and Windows are
planned targets. The final device inventory will be documented before experiments. Where
feasible, the evaluation will compare transfer without target calibration, transfer with a
fixed small calibration budget, and a target-calibrated reference. Configurations used to
evaluate transfer will not be used to fit the initial policy. Calibration cost will be
reported separately from evaluation results. These comparisons characterize transfer between
complete deployment environments and do not identify independent effects of hardware and
operating system.

\section{Scope and intended contributions}
\label{scope-and-intended-contributions}

The study is limited to explicitly registered computations whose inputs and outputs can be
transferred between isolates and whose execution does not directly require the UI. Selection
occurs before task execution. Persistent-worker concurrency and lifecycle will be controlled
explicitly; selecting pool capacity is not itself the proposed contribution. Candidate
workload families include structured-data processing, compression, hashing, indexing and
numerical transformations. UI operations, ordinary network waiting, arbitrary migration of
running code and changes to the Flutter engine or operating-system scheduler are outside the
scope. Flutter Web is excluded because its execution mechanisms differ from the native isolate
model \citep{flutter-isolates}.

The intended contributions are a measured account of execution costs, an open-source Flutter
package exposing the selector through an application-facing interface, and an evaluation under
varying task and UI conditions. Practical evidence about calibration and cases where a fixed
strategy suffices is part of the expected result. Effectiveness and generality will be
established by experiments. The study is intended to help developers choose execution
strategies and understand the costs and limits of application-level adaptation.

\section{Thesis organization}
\label{thesis-organization}

Chapter 2 reviews the execution model, available libraries, runtime selection, UI-aware
adaptation and relevant evaluation methods. The subsequent methodology chapter will define the
artifact, experimental variables, metrics and analysis procedure. Implementation and
evaluation chapters will describe the prototype and report the collected results. The
discussion will relate those results to the research questions and threats to validity,
followed by conclusions and directions for further work.
\chapter{Literature Review}
\label{literature-review}

\section{Review scope and approach}
\label{review-scope-and-approach}

This review examines research relevant to execution-context selection for CPU-bound tasks in
native Flutter applications. It combines official runtime documentation for mechanism
semantics, package documentation and selected source code for engineering capabilities, and
scholarly publications for selection methods, UI-aware adaptation and evaluation. A systematic
review of self-adaptive systems provides field orientation \citep{wong2022}. The present
review groups and compares sources by these questions; it is an exploratory review rather than
a systematic review.

The search was updated on 1 November 2026 using scholarly web discovery, publisher sites,
author pages and institutional repositories. Crossref, a DOI registration agency exposing
publisher-deposited metadata, supported bibliographic checks; OpenAlex, an open scholarly
index, supported forward-citation searches \citep{crossref-api, openalex-docs}. These services
were selected for accessible metadata and citation links, not as indicators of study quality.
Selection favored execution choice, granularity, communication, responsiveness and
calibration. The updated search also covered mobile computation offloading, recent OS-level
interaction-aware scheduling and empirical evidence of main-thread performance problems.
A supplementary search used Flutter/Dart scheduling and isolate terms
independently of the Wagle and iAware citation chains. Its scope, candidate exclusions and
stopping rule are recorded in the companion search log. Detailed comparisons in this review
use inspected full-text sections; inaccessible candidates remain in the working reading queue.

\section{Execution costs and responsiveness}
\label{execution-costs-and-responsiveness}

\subsection{Evidence of main-thread blocking in applications}
\label{evidence-of-main-thread-blocking}

Responsiveness problems have several mechanisms. Android's documentation identifies
excessive computation on the main thread as a cause of unresponsiveness, alongside
blocking I/O and synchronization. Its examples include constructing a complex in-memory
structure and computing a game move \citep{android-responsive}. An input-dispatch ANR
involves a timeout measured in seconds. Frame-scale jank can arise much earlier, so ANR
occurrence and frame-budget exceedances must be evaluated as different outcomes.

Flutter provides a concrete application example in which parsing and converting a large
JSON response causes jank, and moves this computation to a background isolate
\citep{flutter-json}. This supports including structured-data processing among the
experimental workloads. It is an implementation example rather than a measurement of how
frequently deployed Flutter applications suffer from this problem.

Empirical Android research supplies complementary application evidence. Lin et al.
identify rendering, I/O and scheduling causes of poor responsiveness \citep{lin2024}.
Liao et al.~compare performance issues reported in application reviews, Stack Overflow
discussions and GitHub issues and commits with the academic literature. Their classification
includes heavy operations and image decoding on the main thread, as well as rendering,
I/O and synchronization problems \citep[Table~1]{liao2026performance}. The comparison here
uses the inspected author manuscript identified in the reference entry. Its manually
classified, issue-oriented sample establishes concrete problem patterns; it does not estimate
their prevalence across all Android or Flutter applications.

MUSched evaluates foreground interactions in applications alongside a background
picture-in-picture video call \citep[Table~3]{musched2026}. These studies demonstrate the
relevance of resource contention and execution dependencies to application responsiveness.
Three mechanisms must be distinguished: the main thread can be running expensive code,
blocked awaiting a lock or IPC result, or runnable but waiting for CPU time
\citep[Secs. 2--3, 6]{musched2026}. Moving eligible computation to a worker can free the
main isolate to process events, but additional workers can still compete for CPU time.
It does not by itself remove an unrelated lock or IPC dependency. The proposed evaluation
therefore separates deliberately submitted CPU work from background competition and records
both task completion and frame outcomes.

\subsection{UI performance and its limits}
\label{ui-performance-and-its-limits}

The display refresh interval is the time between refresh opportunities: about 16.7 ms at 60 Hz
and 8.3 ms at 120 Hz. Flutter's \texttt{buildDuration} measures the frame-processing interval from
build start to submission for rendering. It includes more than an individual widget build.
\texttt{rasterDuration} measures the raster stage, while \texttt{totalSpan} runs from vsync
start to raster finish. Because stages can overlap across frames, total pipeline latency and
the throughput of frame production are different quantities
\citep{flutter-frame, flutter-build}.

CPU work can also delay the start of frame processing. \texttt{vsyncOverhead} measures the
interval from receiving vsync to build start \citep{flutter-vsync}. A short build duration can
therefore coexist with a long preceding wait. Frame-start intervals under continuous animation
provide a complementary observation of irregular frame production; idle periods must be
excluded, and the display refresh rate must be recorded. Neither these intervals nor pipeline
completion alone establishes the time at which pixels were physically displayed or the full
response to an input event.

Flutter delivers frame-timing reports in batches, approximately once per second in release and
every 100 ms in profile/debug, except for the first report \citep{flutter-timings}. The
selector therefore receives delayed history. Measurement should use profile or release because
debug adds development checks and runtime overhead that can distort the observations
\citep{flutter-frame}.

Lin et al.~distinguish responsiveness failures associated with rendering, I/O and scheduling
\citep[Secs. 2.2, 4.2--4.3]{lin2024}. iAware evaluates reducing contention by moving eligible
service work around interaction periods \citep[Secs. 3--5]{iaware2023}. Together these studies
justify examining resource competition while distinguishing it from expensive rendering or
I/O. The proposed selector acts on eligible application computations, so the evaluation must
identify which kinds of delay actually change.

\subsection{Isolate lifecycle and data movement}
\label{isolate-lifecycle-and-data-movement}

Isolates provide separate mutable application state and event loops, with communication
through messages \citep{flutter-isolates}. This semantic isolation should not be interpreted
as complete physical resource separation: the VM documents shared program structures and a
GC-managed heap within an isolate group \citep{dart-vm}. Background workers can therefore
still compete for device resources. Their presence does not establish an independent CPU or
memory domain.

The official APIs specify cost-relevant behavior rather than portable timing constants.
\texttt{Isolate.spawn} gives approximate base memory overhead on the order of 30 KB, without a
universal startup duration or complete process footprint \citep{dart-spawn}.
\texttt{Isolate.run} starts a fresh isolate and returns the result through
\texttt{Isolate.exit} without copying \citep{dart-run}. Ordinary \texttt{SendPort.send} can
instead copy the mutable objects reachable from a message, while recognized immutable objects
can be shared. Delivery is asynchronous, but sending can require work proportional to the
reachable data on the sender side \citep{dart-send}. Representation and return mechanism
therefore matter alongside byte counts. A closure is a function that can retain variables from
its surrounding scope. Sending it can involve those retained objects, including more state
than expected; an unintended captured resource can make the message unsendable. Task
interfaces should make inputs explicit and avoid accidental capture
\citep{dart-run, dart-spawn}.

\texttt{TransferableTypedData} further separates preparation from transport. Construction
takes time proportional to the data size, whereas sending a prepared transferable takes
constant time and transfers the right to materialize it \citep{dart-transferable}. Similarly,
a shutdown request is not a measurement of completed cleanup: \texttt{Isolate.kill} requests
termination according to its control priority \citep{dart-kill}. A cost comparison should
distinguish preparation, startup, queue waiting, computation, result delivery and observed
termination. It must also allow for overlap between operations instead of assuming that every
component adds sequentially to wall time.

\section{Existing Flutter and Dart approaches}
\label{existing-flutter-and-dart-approaches}

The reviewed ecosystem already supports several execution and lifecycle mechanisms. In
\texttt{worker\_manager} 7.2.9, normal execution uses priority-queued dispatch to available
initialized workers. The selected source also exposes \texttt{executeNow}, which invokes the
callback in the calling isolate, and demand-dependent worker creation and disposal through
\texttt{dynamicSpawning} \citep{worker-manager}. Supporting local execution and reusable
workers is therefore existing functionality. Dynamic worker lifetime is not necessarily one
isolate per task: a worker can serve queued work before becoming idle.

\texttt{isolate\_manager} 6.3.2 provides reusable managers and a one-off \texttt{runFunction}
helper following create, start, compute and stop operations \citep{isolate-manager}. Its
reviewed native route returns results by ordinary port messages, unlike the
terminating-isolate result path of \texttt{Isolate.run}. Ending a worker after the message
does not retroactively remove that message's copying or preparation cost. Squadron supplies
another worker-pool design with configurable capacity and generated interfaces
\citep{squadron}. Worker count, retained state, transport, initialization and queue behavior
must therefore be specified for package-based comparisons.

\texttt{frame\_scheduler} 1.0.0 controls execution timing using a frame-derived performance
zone, priorities and duration estimates to execute, defer or drop callbacks locally
\citep{frame-scheduler}. Its \texttt{fps\_monitor.dart} averages
\texttt{FrameTiming.totalSpan} and inverts the result: this is a latency-derived FPS proxy
rather than measured presentation cadence \citep{frame-fps-source}. Deferred submission can
return before callback completion, so awaiting submission alone does not measure task latency.
Its budget check is an admission heuristic rather than preemption. Comparison must account for
completed, deferred and dropped work, since omitting tasks can improve frame statistics.

Flutter itself provides \texttt{SchedulerBinding.scheduleTask} for short work between frames,
ordered by priority. The documentation recommends tasks of up to about a millisecond; under
the default strategy, priorities below animation do not run while an animation is active
\citep{flutter-schedule-task}. This controls admission on the calling isolate, rather than
relocating computation, and cannot interrupt an already running synchronous function.

Additional packages expose related adaptation. \texttt{cpu\_throttle} 1.0.0 divides iterable
work into slices and yields to the event loop, with optional slice adaptation
\citep{cpu-throttle}. Edge-Veda 2.5.0 calibrates workload budgets and adjusts AI-workload
quality or admission while using persistent workers \citep{edge-veda}. These approaches change
workload structure or resource use rather than selecting among three contexts for an unchanged
registered function.

Amplitude Session Replay 1.0.1 is a particularly close engineering comparison because its
adaptive scheduler combines learned operation costs with frame feedback
\citep{amplitude-replay}. In \texttt{frame\_load\_tracker.dart}, frame pressure uses the
larger of build duration plus vsync overhead and raster duration; total pipeline span is
tracked separately. The tracker smooths observations and checks the age of the latest timing
report, excluding stale retrospective load from the budget estimate. \texttt{estimator.dart}
updates duration estimates from measured samples, reacting faster to cost increases than
decreases under the default smoothing settings. \texttt{capture\_policy.dart} combines
estimated capture cost, frame budget, recent frame health and encoder backlog to admit or
postpone capture. Capture-spacing and freshness rules also apply: a maximum-gap rule can
override the spare-budget condition while encoder-capacity limits remain in force.

The SDK additionally learns separate costs for queued image and paragraph operations in
\texttt{resource\_task\_policy.dart}. Its FIFO drain admits work against the available budget
or, when enabled, allows one sufficiently old task to exceed it to prevent starvation. This is
prior art for combining cost history, UI observations and waiting state at application level.
The decision in these inspected paths is whether and when capture-related work may proceed; it
does not select caller-isolate, fresh-isolate or persistent-worker execution for an unchanged
submitted computation. The thesis therefore needs to evaluate the additional placement
decision and its lifecycle and transfer costs, rather than claim novelty from combining task
estimates with frame feedback. This comparison is based on release-source inspection, not a
measured performance comparison.

The reviewed releases thus provide worker reuse, explicit lifecycle control and several forms
of adaptive admission. Their comparison motivates examining a task-history and UI-feedback
policy for execution-context choice, rather than claiming novelty from worker reuse or
feedback alone.

\section{Runtime execution selection and task granularity}
\label{runtime-execution-selection-and-task-granularity}

\subsection{Mobile computation offloading}
\label{mobile-computation-offloading}

Mobile computation offloading studies when application work should execute on a device or
on external infrastructure. Mora et al.'s recent survey organizes mobile cloud computing
across cloud, fog, cloudlet and edge resources, including application partitioning,
decision granularity and context-dependent offloading \citep{mora2024}. Its treatment of
communication and management overhead provides a framework for asking whether relocation
saves enough computation to repay its additional costs.

MAUI makes this trade-off explicit for methods marked as eligible for remote execution.
Its profiler estimates computation, transferable state and network conditions; its solver
selects local or remote execution to reduce device energy under a latency constraint
\citep[Secs. 3--6]{maui2010}. Eligibility constraints retain device-specific work such as
UI operations locally. MAUI also evaluates profiling strategies through an interactive
video game's frame rate \citep[Sec. 5.2.1]{maui2010}. Here FPS is an evaluation outcome;
it should not be confused with a frame-history signal used by the placement solver.

CloneCloud combines static eligibility analysis with dynamic profiling of computation
and migration costs to partition execution between a phone and its clone
\citep[Secs. 3--4]{clonecloud2011}. It migrates a suspended thread's state and later merges
the result, rather than dispatching an explicitly registered task. Optimized partitions
are selected at application launch for the execution and network environment. In its
evaluation, the smallest workloads remain local because migration overhead outweighs
remote execution savings \citep[Sec. 6]{clonecloud2011}. This is evidence for considering
local execution as a meaningful alternative, rather than assuming that offloading wins.

The present study applies execution-choice reasoning to an entirely on-device setting. It
keeps the registered computation and required result fixed while selecting the main
isolate, a fresh isolate or a persistent worker. Relevant costs include isolate startup,
input preparation, result delivery and queue waiting. Flutter frame feedback introduces an
interactive constraint alongside task cost. Remote servers, network allocation and
application partitioning remain outside the experimental scope. Offloading research already
establishes adaptive placement and overhead-aware decisions; its network and energy models
cannot supply the break-even costs for Dart isolate creation, messaging and worker reuse.

\subsection{Selection as a decision problem}
\label{selection-as-a-decision-problem}

Algorithm-selection research describes choosing an alternative according to instance
characteristics and observed performance \citep{smithmiles2009}. The analogy is useful even
though the proposed modes execute the same computation rather than different solving
algorithms. Task type, data representation, repetition and worker state are candidate
descriptors; the selection criterion must include both task and UI outcomes.

SATzilla makes feature computation part of the selection cost and includes fallback behavior
when feature extraction is unsuitable \citep{satzilla2008}. ASlib distinguishes the strongest
fixed algorithm from a virtual selector choosing the best alternative for each stored instance
\citep{aslib2016}. Together, these works suggest two questions for this thesis: whether the
modes have complementary strengths and whether exploiting them costs less than the potential
benefit. A retrospective per-task winner can provide descriptive context, but interacting
queues and UI effects prevent treating it automatically as an attainable online bound.

Online algorithm selection also has to learn from incomplete observations. Gagliolo and
Schmidhuber discuss training cost, online model updates and exploration versus exploitation
\citep[Secs. 1, 5--6]{gagliolo2006}. Their portfolio allocates computation time among
concurrently run algorithms, whereas this thesis chooses one execution mode before each
submitted task starts. In our setting, executing one mode reveals its cost but not the other
modes' costs on the same submission. A selector that never revisits alternatives can retain
inaccurate estimates, particularly after workload changes. This motivates documenting initial
calibration, controlled probes or estimate aging and counting their cost. Measuring all modes
offline does not make those observations available to the deployed selector; a bandit
algorithm is one possible response, not a requirement of the thesis.

\subsection{The closest execution-choice precedent}
\label{the-closest-execution-choice-precedent}

Wagle et al.~investigate choosing inline execution or a new asynchronous HPX lightweight task
\citep[pp.~4--6]{wagle2019}. Their baseline already uses history: after enough observations,
duration is compared with fixed lower and upper thresholds, retaining the previous mode
between them. The adaptive extension tunes the threshold through APEX and stops after
convergence. Runtime execution choice, history and threshold adaptation are therefore
established parts of the comparison.

Results depend on the evaluated platform: adaptive tuning produces strong gains in some AMD
cases but no average improvement in the reported Intel experiment \citep[p.~8]{wagle2019}.
This supports including unfavorable outcomes and a strong history-based control. The study
uses iterative numerical work on server processors; Dart isolate lifecycle, message transfer
and the UI responsibilities of the main isolate introduce different costs and constraints.

\subsection{Granularity and overhead research}
\label{granularity-and-overhead-research}

Bora et al.~combine queueing analysis, simulation and experiments to study how tiny tasks
balance load distribution against management overhead \citep{bora2023}. The 2023 journal
version evaluates Spark and hybrid Slurm/MPI/OpenMP, extending the earlier Spark-focused
manuscript. It distinguishes useful service, waiting and overhead rather than treating
parallel execution time as a single quantity. This motivates component measurements and
latency distributions. However, splitting a job into many tasks differs from selecting a
context for an unchanged submitted computation, and its fitted runtime constants cannot become
Dart thresholds.

Furmento et al.~(2025) adapt granularity online in StarPU. Their Splitter decides whether to
expand a recursive task before resource scheduling \citep[Sec. 3]{furmento2025}. Periodic
linear programming uses performance estimates and available work; calibrated models and
parameters are still required, and memory transfers remain outside that model
\citep[Secs. 4.2--4.4]{furmento2025}. The dense linear-algebra evaluation includes losing
cases on smaller inputs \citep[Sec. 5.3]{furmento2025}; a separate experiment distinguishes
decision and splitting overhead \citep[Sec. 5.5]{furmento2025}. These findings support
measuring policy cost and current parallelism, while the actions remain task decomposition on
CPU/GPU resources rather than Dart execution-context choice.

\section{UI-aware adaptation and feedback organization}
\label{ui-aware-adaptation-and-feedback-organization}

\subsection{Interaction and resource-aware systems}
\label{interaction-and-resource-aware-systems}

iAware combines touch activity, frame-rate information and resource observations to classify
interaction state \citep[Secs. 3.1--3.2]{iaware2023}. Its model is trained offline, including
separate datasets for device series. Eligible system-service work is advanced or deferred
after dependency checks \citep[Secs. 3.3--3.4]{iaware2023}. Evaluation uses two Android
smartphones and compares Schedutil, SmartSwap and zTT, separating FPS and launch latency and
reporting monitoring overhead \citep[Secs. 4, 5.1, 5.4]{iaware2023}. The implementation
changes service/library code, so its deployment authority differs from an ordinary application
library.

For the present study, iAware establishes that UI-aware adaptation already combines several
kinds of observations. It also demonstrates why task eligibility, model preparation and device
specificity belong in the comparison. Its action is changing when system work occurs. The
proposed selector changes where a registered application computation executes, while retaining
its required result. These distinctions concern decision authority and task semantics, rather
than the presence of a feedback loop alone.

MobiRL optimizes smoothness and power through CPU/GPU frequency limits. The deployed actor
observes the system every ten rendered frames \citep[Secs. 4.3--4.4, 5.4.6]{dou2025}. The
paper separates feature-collection and inference costs, supporting explicit accounting for
observation and decisions \citep[Sec. 5.8]{dou2025}. Privileged Android integration and
possible hardware-dependent retraining differ from deploying a portable Flutter library
\citep[Sec. 6]{dou2025}. Its system-state adaptation is relevant, while its action space and
training cost must remain explicit in the comparison.

A recent OS-level approach is MUSched, which assigns interaction-critical threads to a VIP
scheduling class between real-time and ordinary fair-scheduled work. Framework events
identify critical work, while eBPF-based policies govern its scheduling and bounded priority
boosts \citep[Secs. 4--5]{musched2026}. It propagates interaction priority through lock
and Binder dependencies, addressing cases where a main thread waits for lower-priority
work; the motivating example follows a main-thread IPC request
\citep[Fig.~1]{musched2026}. Its evaluation distinguishes CPU scheduling waits and
dependency-related waits and includes concurrent foreground/background interactions
\citep[Sec. 6, Table~3]{musched2026}. Kernel and framework integration give it authority
over competing threads across the system, beyond an application library's task dispatch.

Zheng et al.~address a different layer: process management. Their measurements associate
slow UI responses with resource competition from background applications, including
applications whose nominal lifecycle state does not reflect their actual user-visible
activity \citep[Secs. 3--4]{zheng2024process}. Their redesigned process-management model
uses runtime behavior to govern background resource retention
\citep[Secs. 5--6]{zheng2024process}. The study was published in 2024, while its
large-scale observational measurements were collected in 2022. Its vendor-specific
population and system intervention should not be read as a current prevalence estimate for
Flutter applications.

These systems distinguish several useful actions: deferring service work, changing frequency
limits, prioritizing dependency chains and managing background processes. Choosing a Dart
execution context is another action with a narrower scope. It can change the computation
occupying the main isolate and the amount of competing worker activity, but cannot grant
OS scheduling priority or control other applications. Their findings motivate measuring
waiting and resource competition alongside the submitted task's execution time.

Lv et al.~take another approach by restructuring rendering jobs according to whether event
results require display \citep{lv2025}. Their study distinguishes frame production from event
response and evaluates the technique using an instrumented synthetic Android application. This
is useful evidence about metric selection and an alternative intervention. A selector for
ordinary computations does not restructure rendering. Selectively skipping rendering jobs
changes the work performed, so comparisons should distinguish completed application events
from produced frames.

\subsection{MAPE-K and policy design}
\label{mape-k-and-policy-design}

Autonomic-computing research established systems that observe and adjust their operation
toward stated goals \citep{kephart2003}. Rainbow separates reusable adaptation infrastructure
from system-specific observations, models and strategies \citep{rainbow2004}. MAPE-K provides
a related organization into monitoring, analysis, planning, execution and shared knowledge;
Arcaini et al.~explicitly model such loops and their properties \citep{arcaini2015}. These
concepts provide vocabulary for the proposed architecture, rather than evidence that any
particular selection policy will be effective.

For this project, monitoring could collect task history, queue state and frame reports;
analysis could estimate pressure or relative costs; planning could select an eligible context;
execution could dispatch work; and knowledge could retain estimates and calibration. This
mapping is a proposed design interpretation. It does not require five separate services or
commit the implementation to a formal model. A named reference architecture alone would not
distinguish the thesis from existing adaptive systems.

Filieri et al.~connect software adaptation to measurable goals, disturbances, transients and
delayed observations \citep{filieri2015}. These concerns apply to stale Flutter reports and
the overhead of changing decisions. MAPE-K organizes the available observations and actions;
the performance of our policy will be tested against fixed execution strategies, a calibrated
hybrid policy with fixed thresholds, and the same selector without UI signals.

\section{Evaluation and calibration implications}
\label{evaluation-and-calibration-implications}

The evaluation must distinguish frame-stage work, waiting before frame processing and
application-task completion. Build/raster distributions and their budget exceedances describe
stage cost; \texttt{vsyncOverhead} and frame-start intervals under continuous animation cover
complementary delays. Task latency includes queue waiting and result delivery, while
throughput counts completed work \citep{flutter-frame, flutter-vsync, iaware2023}. The pilot
should introduce controlled main-isolate blocking and background CPU competition to check
which measures respond. The experiments will record the Flutter SDK version, target
operating system and UI/platform-thread configuration. With merged threads, CPU-bound work
in the main isolate occupies the platform main thread. Native work on that thread is an
additional source of load, so a frame delay must not automatically be attributed to a
submitted task. The selector controls only explicitly registered computations. Input-event response may be added with explicit start/end points; it
must not be labelled physical input-to-display latency without presentation evidence.

The comparison should include a calibrated hybrid policy that selects an execution mode
using task characteristics and worker availability. Its thresholds will be fitted on separate
calibration runs and held fixed during evaluation. This comparator tests whether a practical
rule configured in advance is sufficient. To isolate UI feedback, compare the same
cost/history/queue policy with and without frame signals, keeping initial calibration, non-UI
inputs and learning budgets matched. History can then evolve differently as a consequence of
decisions. Fixed main, fresh-isolate and persistent-worker strategies remain separate
comparators for characterizing execution costs. Controlled policy comparisons will use the
same execution mechanisms, data representations and worker capacity. Record completed work,
latency tails, memory, monitoring and decision costs, including optional probes. Smoothing
frames by dropping required work does not establish a benefit.

Existing worker packages also provide practical reference implementations.
\texttt{worker\_manager} is an initial candidate, with \texttt{isolate\_manager} as an
alternative \citep{worker-manager, isolate-manager}. An integration pilot will establish
which package configuration supports the selected workloads and completion semantics. The
package comparison will use the same task functions and inputs, with worker count,
initialization and lifecycle settings documented. Input preparation, message transfer and
result delivery will be included in end-to-end measurements. Differences in these mechanisms
will be reported as package-level costs rather than attributed solely to the selection
policy. Admission schedulers such as \texttt{frame\_scheduler} require separate treatment
because they may defer or drop work \citep{frame-scheduler}. Any such comparison must account
for completed, deferred and dropped tasks together with task latency and frame outcomes.

Performance-evaluation research also warns against relying on a single favorable run. Georges
et al.~distinguish startup from steady-state behavior and discuss variability and aggregation
\citep{georges2007}. Barrett et al.~show that stable post-warmup behavior cannot simply be
assumed in the JIT runtimes they studied \citep{barrett2017}. These findings motivate
retaining run-level data and examining stability; they do not establish identical behavior in
Flutter's native AOT execution. Cold worker initialization, warm dispatch and policy
adaptation phases should be identifiable in the planned experiment.

Biørn-Hansen et al.'s cross-platform study, including Flutter, demonstrates the value of
explicit comparable operations and hardware diversity \citep{biornhansen2020}. Its historical
framework rankings and platform-feature workloads do not answer the isolate question. For RQ3,
calibration should be evaluated as part of deployment cost. Where the available device
inventory permits, configurations used for initial calibration should be separated from those
used to evaluate transfer. Each configuration combines hardware and an operating system.
Results should characterize transfer across these configurations without attributing observed
differences independently to hardware or operating system. Calibration budgets should report
both the number of task measurements and elapsed calibration time. Evaluation runs must be
separate from calibration runs. Any online updates permitted during evaluation should use
matched budgets across policy variants. Different input representations, retained worker state
or greater pool capacity must not be mistaken for benefits of the selection rule itself.

Calibration also needs a precise meaning. The reviewed approaches include offline training of
a state classifier, fitting runtime performance models and tuning a selection threshold online
\citep{iaware2023, furmento2025, wagle2019}. These preparations have different data
requirements and costs. For this thesis, per-device cost measurement, initial rule fitting and
subsequent online updates should be identified separately. A policy adapting during execution
is not thereby free of prior calibration.

MAUI supplies a concrete example of observation cost: estimating transferable state can
require serialization even when the method eventually runs locally. Its incremental state
transfer also makes later transfer costs depend on earlier placement choices
\citep[Sec. 5.2.1]{maui2010}. Thus, fresh estimates are not free, and the policy can change
the costs it observes. For the proposed selector, feature collection, timing updates and any
calibration probes should be counted alongside dispatch. Reusing estimates may reduce this
cost but requires an explicit rule for aging or refreshing them when workload conditions
change. This supports measuring calibration cost rather than assuming that more frequent
measurement always improves the overall result.

\section{Comparative synthesis and research gap}
\label{comparative-synthesis-and-research-gap}

The reviewed approaches intervene at different levels. Table 2.1 summarizes the decisions most
relevant to the thesis. Observations and objectives describe the cited work's setting, not
equivalent measurements that can be compared numerically across papers.

\begin{longtable}{|
  >{\raggedright\arraybackslash}p{(\linewidth - 8\tabcolsep - 5\arrayrulewidth) * \real{0.2500}}|
  >{\raggedright\arraybackslash}p{(\linewidth - 8\tabcolsep - 5\arrayrulewidth) * \real{0.2500}}|
  >{\raggedright\arraybackslash}p{(\linewidth - 8\tabcolsep - 5\arrayrulewidth) * \real{0.2500}}|
  >{\raggedright\arraybackslash}p{(\linewidth - 8\tabcolsep - 5\arrayrulewidth) * \real{0.2500}}|}
\caption{Decision scope of the closest reviewed approaches.}\tabularnewline
\hline
\begin{minipage}[b]{\linewidth}\raggedright
Approach
\end{minipage} & \begin{minipage}[b]{\linewidth}\raggedright
Observations
\end{minipage} & \begin{minipage}[b]{\linewidth}\raggedright
Decision and objective
\end{minipage} & \begin{minipage}[b]{\linewidth}\raggedright
Preparation and evaluation
\end{minipage} \\
\hline
\endfirsthead
\hline
\begin{minipage}[b]{\linewidth}\raggedright
Approach
\end{minipage} & \begin{minipage}[b]{\linewidth}\raggedright
Observations
\end{minipage} & \begin{minipage}[b]{\linewidth}\raggedright
Decision and objective
\end{minipage} & \begin{minipage}[b]{\linewidth}\raggedright
Preparation and evaluation
\end{minipage} \\
\hline
\endhead
\hline
\endlastfoot
  MAUI / CloneCloud \citep{maui2010, clonecloud2011}
    & Computation, state transfer and network costs
    & Device or remote execution; energy or execution cost
    & Runtime profiling / profiled partitions; phone--server tests
    \\ \hline
%
  Wagle et al. \citep{wagle2019}
    & Per-primitive history
    & Inline or asynchronous; numerical runtime
    & Threshold tuning; HPX, no UI
    \\ \hline
%
  Bora et al. \citep{bora2023}
    & Cost/load distributions
    & Granularity and queueing analysis
    & Fitted overhead model; analysis, simulation and experiments
    \\ \hline
%
  StarPU Splitter \citep{furmento2025}
    & Performance estimates; available work
    & Expand recursive tasks; CPU/GPU throughput
    & Calibrated models; linear-algebra experiments
    \\ \hline
%
  iAware \citep{iaware2023}
    & Touch, FPS, resource state
    & Advance/defer service work
    & Device-series training; Android system changes
    \\ \hline
%
  MobiRL \citep{dou2025}
    & System state; trained actor
    & Frequency limits; smoothness and power
    & Training and retraining; privileged Android tests
    \\ \hline
%
  MUSched \citep{musched2026}
    & Interaction-critical threads; IPC/lock dependencies
    & VIP priority and dependency propagation; responsiveness
    & Kernel/framework integration; mobile application tests
    \\ \hline
%
  Dart packages \citep{worker-manager, isolate-manager, squadron}
    & Queue and worker state; configuration
    & Dispatch and worker lifetime
    & Selected source and docs; not benchmarked here
    \\ \hline
%
  frame\_scheduler \citep{frame-scheduler}
    & Latency proxy; priority; cost estimate
    & Local execution, deferral or dropping
    & Configured estimates; source inspection
    \\ \hline
%
  Built-in / adjacent adaptation \citep{flutter-schedule-task, cpu-throttle, edge-veda}
    & Priorities; slice estimates; domain telemetry
    & Admission, slicing or workload control
    & API and selected source; not benchmarked here
    \\ \hline
%
  Amplitude Session Replay \citep{amplitude-replay}
    & Learned costs; frame pressure; timing age; waiting and backlog
    & Capture and resource admission; no three-context choice in inspected paths
    & Online updates; 1.0.1 source inspection
    \\ \hline
%
  Proposed study
    & Tasks, data and history; queue; frames
    & Main, fresh or persistent; frame and task trade-offs
    & Calibration and transfer (RQ3), not yet evaluated
    \\
%
\end{longtable}

Three conclusions shape the investigation. Runtime execution selection, computation
offloading and interaction-aware adaptation are established. Their decisions span method
placement, task granularity, service timing and OS resource control. The proposed study
therefore cannot claim novelty from adaptive relocation or UI awareness alone. Dart lifecycle,
message transfer and persistent worker state require their own cost characterization.
Adaptation also consumes observation, decision and calibration resources and
learns from selectively observed executions. Finally, frame feedback is delayed and can
describe several different kinds of latency. Its contribution should be tested by disabling
only UI signals in an otherwise matched selector. Comparison with a calibrated hybrid policy
with fixed thresholds will establish whether runtime adaptation adds value beyond a rule
configured in advance. Package-level references will provide a complementary assessment of
practical execution costs.

The research gap is the unresolved empirical relationship between main-isolate execution,
fresh-isolate execution and persistent dispatch costs, task characteristics and UI conditions
in the scoped native Flutter workload. In particular, the reviewed systems do not establish
when delayed frame feedback adds value beyond a cost/history/queue selector for these three
alternatives. RQ1 characterizes the alternatives, RQ2
evaluates context selection and the additional value of UI feedback, and RQ3 examines
calibration budgets and transfer across device--OS configurations. This leads to a controlled
prototype evaluation with matched useful
work, complementary frame observations and explicit accounting for policy overhead.
\clearpage
\addcontentsline{toc}{chapter}{References}
\begingroup
\setlength{\bibsep}{0.35\baselineskip}
\interlinepenalty=10000
% Generated by IEEEtranN.bst, version: 1.14 (2015/08/26)
\begin{thebibliography}{50}
\providecommand{\natexlab}[1]{#1}
\providecommand{\url}[1]{#1}
\csname url@samestyle\endcsname
\providecommand{\newblock}{\relax}
\providecommand{\bibinfo}[2]{#2}
\providecommand{\BIBentrySTDinterwordspacing}{\spaceskip=0pt\relax}
\providecommand{\BIBentryALTinterwordstretchfactor}{4}
\providecommand{\BIBentryALTinterwordspacing}{\spaceskip=\fontdimen2\font plus
\BIBentryALTinterwordstretchfactor\fontdimen3\font minus
  \fontdimen4\font\relax}
\providecommand{\BIBforeignlanguage}[2]{{%
\expandafter\ifx\csname l@#1\endcsname\relax
\typeout{** WARNING: IEEEtranN.bst: No hyphenation pattern has been}%
\typeout{** loaded for the language `#1'. Using the pattern for}%
\typeout{** the default language instead.}%
\else
\language=\csname l@#1\endcsname
\fi
#2}}
\providecommand{\BIBdecl}{\relax}
\BIBdecl

\bibitem[{Flutter team}({\natexlab{a}})]{flutter-isolates}
\BIBentryALTinterwordspacing
{Flutter team}, ``Concurrency and isolates,'' Flutter documentation, accessed:
  Oct. 30, 2026. [Online]. Available:
  \url{https://docs.flutter.dev/perf/isolates}
\BIBentrySTDinterwordspacing

\bibitem{flutter-thread-merge}
Flutter team, ``Merged threads on macOS and Windows,'' Flutter documentation,
accessed: Nov.~1, 2026. [Online]. Available:
\url{https://docs.flutter.dev/release/breaking-changes/macos-windows-merged-threads}

\bibitem[{Flutter team}({\natexlab{b}})]{flutter-frame}
\BIBentryALTinterwordspacing
------, ``{FrameTiming} class,'' Flutter API documentation, accessed: Oct. 30,
  2026. [Online]. Available:
  \url{https://api.flutter.dev/flutter/dart-ui/FrameTiming-class.html}
\BIBentrySTDinterwordspacing

\bibitem[Lin et~al.(2024)Lin, Liu, Li, Qian, Li, Xiong, and Liu]{lin2024}
\BIBentryALTinterwordspacing
H.~Lin, C.~Liu, Z.~Li, F.~Qian, M.~Li, P.~Xiong, and Y.~Liu, ``Aging or
  glitching? what leads to poor {Android} responsiveness and what can we do
  about it?'' \emph{IEEE Transactions on Mobile Computing}, vol.~23, no.~2, pp.
  1521--1533, 2024, accepted author manuscript consulted
  (\url{https://www.linhao.me/pdf/tmc23-responsiveness.pdf}); section locators
  refer to this version. [Online]. Available:
  \url{https://doi.org/10.1109/TMC.2023.3237716}
\BIBentrySTDinterwordspacing

\bibitem[{Dart team}({\natexlab{a}})]{dart-run}
\BIBentryALTinterwordspacing
{Dart team}, ``{Isolate.run},'' Dart API documentation, accessed: Oct. 30,
  2026. [Online]. Available:
  \url{https://api.dart.dev/dart-isolate/Isolate/run.html}
\BIBentrySTDinterwordspacing

\bibitem[{Dart team}({\natexlab{b}})]{dart-send}
\BIBentryALTinterwordspacing
------, ``{SendPort.send},'' Dart API documentation, accessed: Oct. 30, 2026.
  [Online]. Available:
  \url{https://api.dart.dev/dart-isolate/SendPort/send.html}
\BIBentrySTDinterwordspacing

\bibitem[Wagle et~al.(2019)Wagle, Monil, Huck, Malony, Serio, and
  Kaiser]{wagle2019}
\BIBentryALTinterwordspacing
B.~Wagle, M.~A.~H. Monil, K.~Huck, A.~D. Malony, A.~Serio, and H.~Kaiser,
  ``Runtime adaptive task inlining on asynchronous multitasking runtime
  systems,'' in \emph{Proceedings of the 48th International Conference on
  Parallel Processing (ICPP)}, 2019, pp. 1--10, article 76. [Online].
  Available: \url{https://doi.org/10.1145/3337821.3337915}
\BIBentrySTDinterwordspacing

\bibitem[Furmento et~al.(2025)Furmento, Guermouche, Lucas, Morin, Thibault, and
  Wacrenier]{furmento2025}
\BIBentryALTinterwordspacing
N.~Furmento, A.~Guermouche, G.~Lucas, T.~Morin, S.~Thibault, and P.-A.
  Wacrenier, ``Optimizing parallel heterogeneous system efficiency: Dynamic
  task graph adaptation with recursive tasks,'' \emph{Journal of Parallel and
  Distributed Computing}, vol. 205, p. 105157, 2025, author manuscript
  HAL-05199066v1, 2025
  (\url{https://hal.science/hal-05199066/file/paper_unreplaced.pdf}); section
  locators refer to this version. [Online]. Available:
  \url{https://doi.org/10.1016/j.jpdc.2025.105157}
\BIBentrySTDinterwordspacing

\bibitem[Zheng et~al.(2023)Zheng, Li, Xiong, Liu, Ji, Zhu, and Yu]{iaware2023}
\BIBentryALTinterwordspacing
Y.~Zheng, C.~Li, Y.~Xiong, W.~Liu, C.~Ji, Z.~Zhu, and L.~Yu, ``{iAware}:
  Interaction aware task scheduling for reducing resource contention in mobile
  systems,'' \emph{ACM Transactions on Embedded Computing Systems}, vol.~22,
  no.~5s, pp. 145:1--145:24, 2023. [Online]. Available:
  \url{https://doi.org/10.1145/3609391}
\BIBentrySTDinterwordspacing

\bibitem[{Flutter team}({\natexlab{c}})]{flutter-vsync}
\BIBentryALTinterwordspacing
{Flutter team}, ``{FrameTiming.vsyncOverhead},'' Flutter API documentation,
  accessed: Oct. 31, 2026. [Online]. Available:
  \url{https://api.flutter.dev/flutter/dart-ui/FrameTiming/vsyncOverhead.html}
\BIBentrySTDinterwordspacing

\bibitem[{Flutter team}({\natexlab{d}})]{flutter-build}
\BIBentryALTinterwordspacing
------, ``{FrameTiming.buildDuration},'' Flutter API documentation, accessed:
  Oct. 31, 2026. [Online]. Available:
  \url{https://api.flutter.dev/flutter/dart-ui/FrameTiming/buildDuration.html}
\BIBentrySTDinterwordspacing

\bibitem[Wong et~al.(2022)Wong, Wagner, and Treude]{wong2022}
\BIBentryALTinterwordspacing
T.~Wong, M.~Wagner, and C.~Treude, ``Self-adaptive systems: A systematic
  literature review across categories and domains,'' \emph{Information and
  Software Technology}, vol. 148, p. 106934, 2022. [Online]. Available:
  \url{https://doi.org/10.1016/j.infsof.2022.106934}
\BIBentrySTDinterwordspacing

\bibitem[{Crossref}()]{crossref-api}
\BIBentryALTinterwordspacing
{Crossref}, ``{REST API},'' Official metadata-service documentation, accessed:
  Oct. 31, 2026. [Online]. Available:
  \url{https://www.crossref.org/documentation/retrieve-metadata/rest-api/}
\BIBentrySTDinterwordspacing

\bibitem[{OpenAlex}()]{openalex-docs}
\BIBentryALTinterwordspacing
{OpenAlex}, ``{OpenAlex Help Center},'' Official scholarly-index documentation,
  accessed: Oct. 31, 2026. [Online]. Available:
  \url{https://help.openalex.org/}
\BIBentrySTDinterwordspacing

\bibitem{android-responsive}
Android Developers, ``Keep your app responsive,'' Android documentation,
accessed: Nov.~1, 2026. [Online]. Available:
\url{https://developer.android.com/topic/performance/anrs/keep-your-app-responsive}

\bibitem{flutter-json}
Flutter team, ``Parse JSON in the background,'' Flutter cookbook,
accessed: Nov.~1, 2026. [Online]. Available:
\url{https://docs.flutter.dev/cookbook/networking/background-parsing}

\bibitem{liao2026performance}
D.~Liao, S.~Pan, S.~Yang, Y.~Zhao, Z.~Xing, and X.~Sun, ``A comparative study
of Android performance issues in real-world applications and literature,''
\emph{ACM Transactions on Software Engineering and Methodology}, vol.~35, no.~7,
pp.~1--32, 2026, doi: 10.1145/3771932. [Online]. Available:
\url{https://dl.acm.org/doi/10.1145/3771932}.
Consulted author manuscript: arXiv:2407.05090v3, Oct.~2025, 30~pp.,
\url{https://arxiv.org/abs/2407.05090v3}; Table~1 refers to this manuscript.

\bibitem{musched2026}
J.~Xiao et~al., ``Surviving the impossible trinity: Revisiting CPU scheduling
problem on modern COTS mobile devices (Operational Systems),'' in
\emph{20th USENIX Symposium on Operating Systems Design and Implementation
(OSDI 26)}, 2026, pp.~2353--2366. [Online]. Available:
\url{https://www.usenix.org/conference/osdi26/presentation/xiao}

\bibitem[{Flutter team}({\natexlab{e}})]{flutter-timings}
\BIBentryALTinterwordspacing
{Flutter team}, ``{SchedulerBinding.addTimingsCallback},'' Flutter API
  documentation, accessed: Oct. 30, 2026. [Online]. Available:
  \url{https://api.flutter.dev/flutter/scheduler/SchedulerBinding/addTimingsCallback.html}
\BIBentrySTDinterwordspacing

\bibitem[{Dart team}({\natexlab{c}})]{dart-vm}
\BIBentryALTinterwordspacing
{Dart team}, ``Introduction to {Dart VM},'' Dart SDK runtime documentation, sDK
  commit 416e33b0ffaef062a60ec966bb8c82ed2b83f128; accessed: Oct. 30, 2026.
  [Online]. Available:
  \url{https://github.com/dart-lang/sdk/blob/416e33b0ffaef062a60ec966bb8c82ed2b83f128/runtime/docs/README.md}
\BIBentrySTDinterwordspacing

\bibitem[{Dart team}({\natexlab{d}})]{dart-spawn}
\BIBentryALTinterwordspacing
------, ``{Isolate.spawn},'' Dart API documentation, accessed: Oct. 30, 2026.
  [Online]. Available:
  \url{https://api.dart.dev/dart-isolate/Isolate/spawn.html}
\BIBentrySTDinterwordspacing

\bibitem[{Dart team}({\natexlab{e}})]{dart-transferable}
\BIBentryALTinterwordspacing
------, ``{TransferableTypedData} class,'' Dart API documentation, accessed:
  Oct. 30, 2026. [Online]. Available:
  \url{https://api.dart.dev/dart-isolate/TransferableTypedData-class.html}
\BIBentrySTDinterwordspacing

\bibitem[{Dart team}({\natexlab{f}})]{dart-kill}
\BIBentryALTinterwordspacing
------, ``{Isolate.kill},'' Dart API documentation, accessed: Oct. 30, 2026.
  [Online]. Available:
  \url{https://api.dart.dev/dart-isolate/Isolate/kill.html}
\BIBentrySTDinterwordspacing

\bibitem[{worker\_manager contributors}()]{worker-manager}
\BIBentryALTinterwordspacing
{worker\_manager contributors}, ``{worker\_manager}: release 7.2.9,'' Package
  documentation and selected native source, selected files matched against the
  7.2.9 release archive, not a verified whole-release commit; accessed: Oct.
  30, 2026. [Online]. Available:
  \url{https://github.com/dsrenesanse/worker_manager/blob/5c70ecc12f459b97c476ce3dd7a2da4218d71121/lib/src/scheduling/executor.dart}
\BIBentrySTDinterwordspacing

\bibitem[{isolate\_manager contributors}()]{isolate-manager}
\BIBentryALTinterwordspacing
{isolate\_manager contributors}, ``{isolate\_manager}: release 6.3.2,'' Package
  documentation and selected native source, source tag v6.3.2, commit
  1437cc927af3eaa4910cc62090b4d5c4004c86b8; accessed: Oct. 30, 2026. [Online].
  Available:
  \url{https://github.com/lamnhan066/isolate_manager/blob/1437cc927af3eaa4910cc62090b4d5c4004c86b8/lib/src/isolate_manager.dart}
\BIBentrySTDinterwordspacing

\bibitem[{Squadron contributors}()]{squadron}
\BIBentryALTinterwordspacing
{Squadron contributors}, ``{Squadron}: release 7.4.4,'' Package and
  concurrency-settings documentation, accessed: Oct. 30, 2026. [Online].
  Available: \url{https://pub.dev/packages/squadron/versions/7.4.4}
\BIBentrySTDinterwordspacing

\bibitem[{frame\_scheduler contributors}({\natexlab{a}})]{frame-scheduler}
\BIBentryALTinterwordspacing
{frame\_scheduler contributors}, ``{frame\_scheduler}: release 1.0.0,'' Package
  documentation and selected source, source commit
  c40bac847823579376734d9e7a9e118070f9be83; accessed: Oct. 30, 2026. [Online].
  Available:
  \url{https://github.com/Brah-Timo/frame_scheduler/blob/c40bac847823579376734d9e7a9e118070f9be83/lib/src/scheduler_controller.dart}
\BIBentrySTDinterwordspacing

\bibitem[{frame\_scheduler contributors}({\natexlab{b}})]{frame-fps-source}
\BIBentryALTinterwordspacing
------, ``{frame\_scheduler} 1.0.0: {fps\_monitor.dart},'' Selected release
  source, lines 189--231, accessed: Oct. 31, 2026. [Online]. Available:
  \url{https://github.com/Brah-Timo/frame_scheduler/blob/c40bac847823579376734d9e7a9e118070f9be83/lib/src/fps_monitor.dart}
\BIBentrySTDinterwordspacing

\bibitem[{Flutter team}({\natexlab{f}})]{flutter-schedule-task}
\BIBentryALTinterwordspacing
{Flutter team}, ``{SchedulerBinding.scheduleTask},'' Flutter API documentation,
  accessed: Oct. 31, 2026. [Online]. Available:
  \url{https://api.flutter.dev/flutter/scheduler/SchedulerBinding/scheduleTask.html}
\BIBentrySTDinterwordspacing

\bibitem[{cpu\_throttle contributors}()]{cpu-throttle}
\BIBentryALTinterwordspacing
{cpu\_throttle contributors}, ``{cpu\_throttle}: release 1.0.0,'' Package
  documentation and selected release source, accessed: Oct. 31, 2026. [Online].
  Available: \url{https://pub.dev/packages/cpu_throttle/versions/1.0.0}
\BIBentrySTDinterwordspacing

\bibitem[{Edge-Veda contributors}()]{edge-veda}
\BIBentryALTinterwordspacing
{Edge-Veda contributors}, ``{Edge-Veda}: release 2.5.0,'' Package documentation
  and selected scheduler source, accessed: Oct. 31, 2026. [Online]. Available:
  \url{https://pub.dev/packages/edge_veda/versions/2.5.0}
\BIBentrySTDinterwordspacing

\bibitem[{Amplitude}()]{amplitude-replay}
\BIBentryALTinterwordspacing
{Amplitude}, ``{Amplitude Session Replay Flutter SDK}: release 1.0.1,'' Package
  release and selected capture-admission source, accessed: Oct. 31, 2026.
  [Online]. Available:
  \url{https://pub.dev/packages/amplitude_session_replay/versions/1.0.1}
\BIBentrySTDinterwordspacing

\bibitem{mora2024}
H.~Mora, F.~A. Pujol, T.~Ramirez-Gordillo, and A.~Jimeno, ``Mobile cloud
computing paradigm: A survey of operational concerns, challenges and open
issues,'' \emph{Transactions on Emerging Telecommunications Technologies},
vol.~35, no.~12, article e70020, 2024, doi: 10.1002/ett.70020.
[Online]. Available: \url{https://doi.org/10.1002/ett.70020}

\bibitem{maui2010}
E.~Cuervo et~al., ``MAUI: Making smartphones last longer with code offload,'' in
\emph{Proceedings of the 8th International Conference on Mobile Systems,
Applications, and Services (MobiSys)}, 2010, pp.~49--62,
doi: 10.1145/1814433.1814441. [Online]. Available:
\url{https://doi.org/10.1145/1814433.1814441}

\bibitem{clonecloud2011}
B.-G.~Chun, S.~Ihm, P.~Maniatis, M.~Naik, and A.~Patti, ``CloneCloud: Elastic
execution between mobile device and cloud,'' in \emph{Proceedings of the Sixth
Conference on Computer Systems (EuroSys)}, 2011, pp.~301--314,
doi: 10.1145/1966445.1966473. [Online]. Available:
\url{https://dl.acm.org/doi/10.1145/1966445.1966473}

\bibitem[Smith-Miles(2009)]{smithmiles2009}
\BIBentryALTinterwordspacing
K.~A. Smith-Miles, ``Cross-disciplinary perspectives on meta-learning for
  algorithm selection,'' \emph{ACM Computing Surveys}, vol.~41, no.~1, pp.
  6:1--6:25, 2009. [Online]. Available:
  \url{https://doi.org/10.1145/1456650.1456656}
\BIBentrySTDinterwordspacing

\bibitem[Xu et~al.(2008)Xu, Hutter, Hoos, and Leyton-Brown]{satzilla2008}
\BIBentryALTinterwordspacing
L.~Xu, F.~Hutter, H.~H. Hoos, and K.~Leyton-Brown, ``{SATzilla}:
  Portfolio-based algorithm selection for {SAT},'' \emph{Journal of Artificial
  Intelligence Research}, vol.~32, pp. 565--606, 2008. [Online]. Available:
  \url{https://doi.org/10.1613/jair.2490}
\BIBentrySTDinterwordspacing

\bibitem[Bischl et~al.(2016)Bischl, Kerschke, Kotthoff, Lindauer, Malitsky,
  Fr\'{e}chette, Hoos, Hutter, Leyton-Brown, Tierney, and
  Vanschoren]{aslib2016}
\BIBentryALTinterwordspacing
B.~Bischl, P.~Kerschke, L.~Kotthoff, M.~Lindauer, Y.~Malitsky,
  A.~Fr\'{e}chette, H.~Hoos, F.~Hutter, K.~Leyton-Brown, K.~Tierney, and
  J.~Vanschoren, ``{ASlib}: A benchmark library for algorithm selection,''
  \emph{Artificial Intelligence}, vol. 237, pp. 41--58, 2016. [Online].
  Available: \url{https://doi.org/10.1016/j.artint.2016.04.003}
\BIBentrySTDinterwordspacing

\bibitem[Gagliolo and Schmidhuber(2006)]{gagliolo2006}
\BIBentryALTinterwordspacing
M.~Gagliolo and J.~Schmidhuber, ``Learning dynamic algorithm portfolios,''
  \emph{Annals of Mathematics and Artificial Intelligence}, vol.~47, pp.
  295--328, 2006, issue year 2006; online 2007. Section locators refer to the
  author manuscript at
  \url{https://sferics.idsia.ch/pub/juergen/gagliolo-amai837.pdf}. [Online].
  Available: \url{https://doi.org/10.1007/s10472-006-9036-z}
\BIBentrySTDinterwordspacing

\bibitem[Bora et~al.(2023)Bora, Walker, and Fidler]{bora2023}
\BIBentryALTinterwordspacing
S.~Bora, B.~Walker, and M.~Fidler, ``The tiny-tasks granularity trade-off:
  Balancing overhead versus performance in parallel systems,'' \emph{IEEE
  Transactions on Parallel and Distributed Systems}, vol.~34, no.~4, pp.
  1128--1144, 2023. [Online]. Available:
  \url{https://doi.org/10.1109/TPDS.2022.3233712}
\BIBentrySTDinterwordspacing

\bibitem[Dou et~al.(2025)Dou, Liu, and Xiao]{dou2025}
\BIBentryALTinterwordspacing
X.~Dou, L.~Liu, and L.~Xiao, ``An intelligent scheduling approach on mobile
  {OS} for optimizing {UI} smoothness and power,'' \emph{ACM Transactions on
  Architecture and Code Optimization}, vol.~22, no.~1, pp. 12:1--12:27, 2025.
  [Online]. Available: \url{https://doi.org/10.1145/3674910}
\BIBentrySTDinterwordspacing

\bibitem{zheng2024process}
J.~Zheng et~al., ``Rethinking process management for interactive mobile systems,''
in \emph{Proceedings of the 30th Annual International Conference on Mobile
Computing and Networking (MobiCom)}, 2024, pp.~215--229,
doi: 10.1145/3636534.3649357. [Online]. Available:
\url{https://dl.acm.org/doi/10.1145/3636534.3649357}

\bibitem[Lv et~al.(2025)Lv, Hu, Cui, Yang, Zhou, Deng, and Guan]{lv2025}
\BIBentryALTinterwordspacing
M.~Lv, T.~Hu, M.~Cui, T.~Yang, Y.~Zhou, Q.~Deng, and N.~Guan, ``Improving {UI}
  responsiveness in {Android} by restructured rendering,'' \emph{Journal of
  Systems Architecture}, vol. 168, p. 103580, 2025. [Online]. Available:
  \url{https://doi.org/10.1016/j.sysarc.2025.103580}
\BIBentrySTDinterwordspacing

\bibitem[Kephart and Chess(2003)]{kephart2003}
\BIBentryALTinterwordspacing
J.~O. Kephart and D.~M. Chess, ``The vision of autonomic computing,''
  \emph{Computer}, vol.~36, no.~1, pp. 41--50, 2003. [Online]. Available:
  \url{https://doi.org/10.1109/MC.2003.1160055}
\BIBentrySTDinterwordspacing

\bibitem[Garlan et~al.(2004)Garlan, Cheng, Huang, Schmerl, and
  Steenkiste]{rainbow2004}
\BIBentryALTinterwordspacing
D.~Garlan, S.-W. Cheng, A.-C. Huang, B.~Schmerl, and P.~Steenkiste, ``Rainbow:
  Architecture-based self-adaptation with reusable infrastructure,''
  \emph{Computer}, vol.~37, no.~10, pp. 46--54, 2004. [Online]. Available:
  \url{https://doi.org/10.1109/MC.2004.175}
\BIBentrySTDinterwordspacing

\bibitem[Arcaini et~al.(2015)Arcaini, Riccobene, and Scandurra]{arcaini2015}
\BIBentryALTinterwordspacing
P.~Arcaini, E.~Riccobene, and P.~Scandurra, ``Modeling and analyzing {MAPE-K}
  feedback loops for self-adaptation,'' in \emph{IEEE/ACM 10th International
  Symposium on Software Engineering for Adaptive and Self-Managing Systems
  (SEAMS)}, 2015, pp. 13--23. [Online]. Available:
  \url{https://doi.org/10.1109/SEAMS.2015.10}
\BIBentrySTDinterwordspacing

\bibitem[Filieri et~al.(2015)Filieri, Maggio, Angelopoulos, D'Ippolito,
  Gerostathopoulos, Hempel, Hoffmann, Jamshidi, Kalyvianaki, Klein, Krikava,
  Misailovic, Papadopoulos, Ray, Sharifloo, Shevtsov, Ujma, and
  Vogel]{filieri2015}
\BIBentryALTinterwordspacing
A.~Filieri, M.~Maggio, K.~Angelopoulos, N.~D'Ippolito, I.~Gerostathopoulos,
  A.~B. Hempel, H.~Hoffmann, P.~Jamshidi, E.~Kalyvianaki, C.~Klein, F.~Krikava,
  S.~Misailovic, A.~V. Papadopoulos, S.~Ray, A.~M. Sharifloo, S.~Shevtsov,
  M.~Ujma, and T.~Vogel, ``Software engineering meets control theory,'' in
  \emph{IEEE/ACM 10th International Symposium on Software Engineering for
  Adaptive and Self-Managing Systems (SEAMS)}, 2015, pp. 71--82. [Online].
  Available: \url{https://doi.org/10.1109/SEAMS.2015.12}
\BIBentrySTDinterwordspacing

\bibitem[Georges et~al.(2007)Georges, Buytaert, and Eeckhout]{georges2007}
\BIBentryALTinterwordspacing
A.~Georges, D.~Buytaert, and L.~Eeckhout, ``Statistically rigorous {Java}
  performance evaluation,'' in \emph{Proceedings of the 22nd ACM SIGPLAN
  Conference on Object-Oriented Programming Systems and Applications (OOPSLA)},
  2007, pp. 57--76. [Online]. Available:
  \url{https://doi.org/10.1145/1297027.1297033}
\BIBentrySTDinterwordspacing

\bibitem[Barrett et~al.(2017)Barrett, Bolz-Tereick, Killick, Mount, and
  Tratt]{barrett2017}
\BIBentryALTinterwordspacing
E.~Barrett, C.~F. Bolz-Tereick, R.~Killick, S.~Mount, and L.~Tratt, ``Virtual
  machine warmup blows hot and cold,'' \emph{Proceedings of the ACM on
  Programming Languages}, vol.~1, no. OOPSLA, pp. 52:1--52:27, 2017. [Online].
  Available: \url{https://doi.org/10.1145/3133876}
\BIBentrySTDinterwordspacing

\bibitem[Bi{\o}rn-Hansen et~al.(2020)Bi{\o}rn-Hansen, Rieger, Gr{\o}nli,
  Majchrzak, and Ghinea]{biornhansen2020}
\BIBentryALTinterwordspacing
A.~Bi{\o}rn-Hansen, C.~Rieger, T.-M. Gr{\o}nli, T.~A. Majchrzak, and G.~Ghinea,
  ``An empirical investigation of performance overhead in cross-platform mobile
  development frameworks,'' \emph{Empirical Software Engineering}, vol.~25, pp.
  2997--3040, 2020. [Online]. Available:
  \url{https://doi.org/10.1007/s10664-020-09827-6}
\BIBentrySTDinterwordspacing

\end{thebibliography}
\endgroup
\end{document}
